\documentclass[a4paper, titlepage]{article}

\usepackage[utf8]{inputenc} 
\usepackage[T1]{fontenc}      
\usepackage[francais]{babel}
\usepackage{amsmath, amsthm, amssymb}
\usepackage{tikz}
\usepackage{algorithm, algorithmic}
\usepackage{array}
\usepackage{cite}
\usepackage{tikz}
%\usepackage{hyperref}
\allowdisplaybreaks

\newtheorem{theo}{Théorème}[section]
\newtheorem{lemm}[theo]{Lemme}
\newtheorem{prop}[theo]{Proposition}
\newtheorem{coro}[theo]{Corollaire}
\theoremstyle{definition}
\newtheorem{defi}[theo]{Définition}
\theoremstyle{remark}
\newtheorem{rema}[theo]{Remarque}
\newtheorem{exem}[theo]{Exemple}
\newtheorem{appl}[theo]{Application}

\def\N{\mathbb N}
\def\A{\mathbb A}
\def\Z{\mathbb Z}
\def\Q{\mathbb Q}
\def\R{\mathbb R}
\def\C{\mathbb C}
\def\K{\mathbb K}
\def\F{\mathbb F}
\def\O{\mathcal O}
\def\o{o}
\def\gf{\operatorname{GF}}
\def\frob{\operatorname{Frob}}
\def\card{\operatorname{Card}}
\def\car{\operatorname{car}}
\def\pgcd{\operatorname{pgcd}}
\def\id{\operatorname{id}}
\def\aut{\operatorname{Aut}}
\def\hom{\operatorname{Hom}}
\def\isom{\operatorname{Isom}}
\def\gal{\operatorname{Gal}}
\def\mbf#1{\mathbf{#1}}
\def\NP{\mathbb{NP}}
\def\gen #1{\left\langle#1\right\rangle}
\def\ceil #1{\left\lceil#1\right\rceil}
\def\floor #1{\left\lfloor#1\right\rfloor}

\floatname{algorithm}{Algorithme}
\renewcommand{\algorithmicrequire}{\textbf{Entrée :}}
\renewcommand{\algorithmicensure}{\textbf{Sortie :}}
\renewcommand{\algorithmicend}{\textbf{fin}}
\renewcommand{\algorithmicif}{\textbf{si}}
\renewcommand{\algorithmicthen}{\textbf{alors}}
\renewcommand{\algorithmicelse}{\textbf{sinon}}
\renewcommand{\algorithmicfor}{\textbf{pour}}
\renewcommand{\algorithmicforall}{\textbf{pour tout}}
\renewcommand{\algorithmicdo}{\textbf{faire}}
\renewcommand{\algorithmicwhile}{\textbf{tant que}}
\renewcommand{\algorithmicloop}{\textbf{boucle}}
\renewcommand{\algorithmicrepeat}{\textbf{repéter}}
\renewcommand{\algorithmicuntil}{\textbf{jusqu'à}}
\renewcommand{\algorithmicprint}{\textbf{afficher}}
\renewcommand{\algorithmicreturn}{\textbf{retourner}}
\renewcommand{\algorithmictrue}{\textbf{vrai}}
\renewcommand{\algorithmicfalse}{\textbf{faux}}

\title{\'Etude du cryptosystème de Chor-Rivest}    
\author{Rémi {\sc Clarisse}}        
\date{Mai-Août 2017}       

\begin{document}
\thispagestyle{empty}
\centerline{\includegraphics[width=8cm,height=25mm]{logo-bordeaux.png} \hspace{1cm} \includegraphics[width=8cm,height=25mm]{logo-inria.png}}

\vspace{1cm}

\centerline{\large\bfseries Master 1 de Cryptologie et Sécurité Informatique}
\centerline{\large\bfseries \'Equipe-projet GRACE}
\centerline{\rule{5cm}{2pt}}
\centerline{\bfseries Stage estival 2017}


\begin{center}
\vspace{4cm}

{\Huge\bfseries Rapport de stage \\}
\vspace{0.15cm}\centerline{\rule{10cm}{0.5pt}}\vspace{0.15cm}
{\LARGE\bfseries \'Etude du cryptosystème de Chor-Rivest}
 \vspace{2cm}

{\large Rémi {\sc Clarisse}}

\vspace{6cm}

{\large Tuteurs: Daniel {\sc Augot} et Luca {\sc De Feo}}
\end{center}
\newpage

\section*{Introduction}

Dans l'article~\cite{chorRivest1988}

\newpage
\bibliographystyle{abbrv}
\bibliography{mybib}

\end{document}